\documentclass[11pt,a4paper]{article}
\usepackage[margin=2.5cm,headheight=14pt]{geometry}
\usepackage{amsmath}
\usepackage{siunitx}
\usepackage{booktabs}
\usepackage{fancyhdr}
\usepackage{lastpage}
\usepackage{xcolor}
\usepackage{tabularx}
\usepackage{caption}
\usepackage{microtype}
\usepackage[hidelinks]{hyperref}

\definecolor{band}{RGB}{30,70,110}
\captionsetup{font=small,labelfont=bf}

% Report details used on the cover block and in the header.
\newcommand{\reportnumber}{TR-2026-014}
\newcommand{\reporttitle}{Tensile Testing of 3D-Printed PLA Specimens}

\pagestyle{fancy}
\fancyhf{}
\lhead{\small\reportnumber}
\rhead{\small\reporttitle}
\cfoot{\small Page \thepage\ of \pageref{LastPage}}

\begin{document}

\thispagestyle{empty}
{\color{band}\rule{\linewidth}{4pt}}
\vspace{0.5em}

{\noindent\Huge\bfseries\reporttitle\par}
\vspace{0.4em}
{\noindent\large Technical Report \reportnumber\par}
\vspace{1.2em}

\noindent\begin{tabularx}{\linewidth}{@{}lX@{}}
\textbf{Prepared by} & Lab Member Name, Mechanical Testing Laboratory \\
\textbf{Supervisor} & Supervisor Name \\
\textbf{Date of tests} & 2 to 6 March 2026 \\
\textbf{Status} & Draft for internal review \\
\end{tabularx}

\vspace{1em}
{\color{band}\rule{\linewidth}{0.6pt}}

\section*{Summary}
In three or four sentences, say what was tested, how, and the most important number.

\tableofcontents

\section{Objective}
State what question the tests were designed to answer.

\section{Equipment and materials}
\begin{itemize}
  \item Universal testing machine, \SI{10}{\kilo\newton} load cell
  \item PLA filament, \SI{1.75}{\milli\metre}, printed at \SI{210}{\celsius}
  \item Digital calipers with \SI{0.01}{\milli\metre} resolution
\end{itemize}

\section{Procedure}
\begin{enumerate}
  \item Condition specimens for \SI{48}{\hour} at \SI{23}{\celsius}.
  \item Measure gauge width and thickness at three points.
  \item Load at a crosshead speed of \SI{5}{\milli\metre\per\minute} until fracture.
\end{enumerate}
Ultimate tensile strength is computed as $\sigma_u = F_{\max}/A_0$.

\section{Results}
Table~\ref{tab:results} lists the measured values and Figure~\ref{fig:curve} shows where the stress and strain plot belongs.

\begin{table}[ht]
\centering
\caption{Tensile results by infill orientation (example data).}
\label{tab:results}
\begin{tabular}{l S[table-format=1.0] S[table-format=2.1] S[table-format=1.2] S[table-format=1.1]}
\toprule
{Orientation} & {$n$} & {$\sigma_u$ (\si{\mega\pascal})} & {$E$ (\si{\giga\pascal})} & {Strain at break (\%)} \\
\midrule
0\textdegree  & 5 & 52.3 & 3.41 & 3.2 \\
45\textdegree & 5 & 44.8 & 3.12 & 2.7 \\
90\textdegree & 5 & 31.6 & 2.95 & 1.9 \\
\bottomrule
\end{tabular}
\end{table}

\begin{figure}[ht]
\centering
\fbox{\parbox[c][5.5cm][c]{0.8\linewidth}{\centering\textsf{Figure placeholder}\\[0.3em]\small Replace with a stress and strain plot}}
\caption{Representative stress and strain curves.}
\label{fig:curve}
\end{figure}

\section{Discussion}
Explain trends, likely sources of error and whether the objective was met.

\section{Conclusions and recommendations}
\begin{enumerate}
  \item First conclusion, backed by a number from the results.
  \item Recommendation for design or further testing.
\end{enumerate}

\appendix
\section{Raw measurements}
Attach full data tables or calibration certificates here.

\end{document}
